
% Partie 2 : Format des fichiers vecteurs
\part{Format des fichiers vecteurs}
\label{part:formvec}
\section{Présentation générale}
  Les \emph{fichiers-vecteurs} sont destinés au stockage de séries
temporelles de paramètres mesurés ou calculés. Ils associent
étroitement la donnée (dans un format binaire,
copie directe de la représentation en mémoire de
l'ordinateur), et sa description (dans une en-tête au format ASCII).
%Ce format d

  L'en-tête descriptive contient toutes les informations
utiles (format, unité, appareil ou algorithme d'origine...), \textbf{sauf}
la datation, car celle-ci
est parfois très complexe à gérer (variables
asynchrones). A chaque fichier de données, on associe donc un fichier
temps qui indique l'instant de la mesure de chaque
échantillon. Il y a strictement concordance entre ces deux fichiers :
la date de la i\up{ème} donnée est la i\up{ème}
valeur du fichier temps.

      Naturellement, plusieurs fichiers de données
synchrones sont associés au même fichier \og temps\fg{} : on
ne duplique pas inutilement ces fichiers (dans le cas simplifié
où toutes les variables seraient acquises à la même
fréquence par le même système, il n'y aurait qu'un seul vecteur-temps).

\section{Contenu}

Ces fichiers sont organisés en 2 parties :
\begin{enumerate}
  \item Un en-tête descriptif (souvent nommé \og header\fg{})
  \item Les données en binaire.
\end{enumerate}

\subsection{Format de l'en-tête.}
L'en-tête est un bloc d'une taille multiple de 512
octets, qui contient quelques informations techniques et le descriptif de la
variable dans un bloc ASCII. Ce bloc ASCII est encadré entre deux
marqueurs de début et de fin.
Voici la structure précise de l'en-tête :
\begin{itemize}
  \item Une chaîne de repérage : \og VECT-DATA\texttt{\textbackslash{}n}\fg{}
  \item La taille de l'en-tête (complète, avec les marqueurs, le bloc
ASCII, le remplissage de fin), en octets. Cette taille est forcément un
multiple de 512, et elle est écrite en ASCII cadré à
droite sur 10 caractères suivis d'un saut de ligne 
(\og \texttt{\textbackslash{}n}\fg{}).
  \item le marqueur de début du bloc ASCII, qui se compose 
    d'un caractère \og STX\fg{} (\texttt{\textbackslash{}002} ou \^{}B)
    et d'un caractère \og LF\fg{} (\texttt{\textbackslash{}n}).
  \item Le bloc ASCII, portion de texte documentant la variable, dont la
description précise est donnée plus loin.
  \item Le marqueur de fin du bloc ASCII, qui se compose de :
  \begin{itemize}
    \item une chaîne \og FIN\_HDR\texttt{\textbackslash{}{}n}\fg{}
    \item la suite des caractères \og NP\fg{} (\texttt{\textbackslash{}}0cx ou \^{}L),
       \og ETX\fg{} (\texttt{\textbackslash{}}003 ou \^{}C),
      \og SUB\fg{} (\texttt{\textbackslash{}}1ax ou \^{}Z)
  \end{itemize}
\end{itemize}
Ensuite, cet en-tête, jusqu'à la fin (c.a.d jusqu'au prochain
multiple de 512 octets), est rempli de caractères \og \fg{} (espace).


\subsubsection{Syntaxe du \og bloc-ASCII\fg{}}
Le texte décrivant la variable prend la forme suivante :
\begin{itemize}
  \item une première ligne contenant le nom de la variable (normalement
identique au \texttt{basename} du nom de fichier), encadré de crochets.
  \item les lignes suivantes sous la forme \og \textsf{attribut = valeur}\fg{}
\end{itemize}
Les caractères sont codés selon la norme ASCII, ou
éventuellement selon son extension ISO-8859-1 (ou ISO-Latin1).
Le caractère de fin de ligne est le \og NL\fg{}
(Il n'y a pas actuellement de spécification de longueur de ligne maximale).

%%%\paragraph{attributs obligatoires}
Les attributs suivants sont obligatoires :

\begin{description}
  \item[Nom du paragraphe] (entre crochet) nom du paramètre
et \texttt{basename} du nom de fichier.
  \item[vol] nom unique attribué à la série de mesure. C'est une chaîne de
6 à 8 caractère au format \texttt{PPaann[nn]}, et qui est décrite de façon
détaillée section \ref{sec:AcqId} page \pageref{sec:AcqId}. Pour des raisons
de compatibilité historique, ce champ s'appelle exactement «vol», même pour des
mesures de station de surface, ou à bord de bateaux.
  \item[format] le codage informatique des données. La liste des
  formats connus par la version actuelle d'AIDA est la suivante :

  \begin{description}
    \item[\textsf{Byte}, \textsf{Int\_16}, \textsf{Int\_32}, \textsf{Int\_64}] :\\
       entiers signés
    \item[\textsf{U\_Int\_8}, \textsf{U\_Int\_16}, \textsf{U\_Int\_32}] : \\
       entiers non signés
    \item[\textsf{Ieeefloat\_32}, \textsf{Ieeefloat\_64}] : \\
       Nombres \og flottants\fg{} codés selon la norme ANSI/IEEE 754-1985.
    \item[\textsf{Quanta\_12\_bits},
          \textsf{Quanta\_14\_bits},
          \textsf{Quanta\_16\_bits},
          \textsf{Quanta\_32\_bits}] : \\
       «mots-machine» provenant d'un convertisseur analogique-numérique. \\
       Les \textsf{Quanta\_12\_bits} et \textsf{Quanta\_14\_bits} sont
       stockés dans des mots de 16 bits, avec seulement 12 (ou 14) bits
       significatifs.
    \item[\textsf{Quanta\_14\_bits\_signe}] : \\
       la version signée d'un \textsf{Quanta\_14\_bits}.
    \item[\textsf{Arinc\_SY429}] : \\
       codage selon la norme aéronautique ARINC-SY429 (mot de 64 bits)
    \item[\textsf{BC635\_Time}] : Codage produits par la carte BC635.
    \item[\textsf{Special}] : Codages particuliers.
    \item[\textsf{Hpfloat\_32}] : \\
       nombre «flottantsi» codés à la manières des machines HP-1000.
    \item[\textsf{NCAR2D\_time}] : Le champ de datation des échantillons 
«2D» enregistrés par la carte développée au NCAR.
  \end{description}

  Dans les noms de formats, la casse (majuscule/minuscule) n'est pas signifiante.
  \item[endianness] indique l'ordre de la séquences d'octets des données. Ce champ peut prendre la valeur \texttt{little} ou \texttt{big}. Ce champ a été ajouté en 2021. Afin de conserver la compatibilité avec les anciens vecteurs, si ce champ n'est pas présent alors on considère endianness = big. 
  \item[dimension] taille en octet d'une
représentation de la variable. Si la valeur de l'attribut
\textsf{format} était un type connu du
langage C, on aurait : \textsf {dimension = sizeof(format)}. En
général, donc, \textsf{dimension} peut se
déduire de \textsf{format}, l'exception étant bien sûr le format 
\textsf{Special}.
  \item[dim\_array]  Dans le cas où chaque échantillon de la
variable est représenté par un tableau au lieu d'une simple
valeur (par exemple les spectres des FSSP), on indique ici les dimensions de
ce tableau. Dans le cas où le tableau a plus d'une dimension (une
image, par exemple), les tailles sont indiquées séparées par un \textsf{|}.
L'ordre des tailles indiquées suit la convention
du langage C. Ce champ n'est pas obligatoire; en son absence on suppose
que la variables est scalaire (un seul échantillon par instant de mesure).
   \item[ref\_time] Nom (\og logique\fg{}, sans le \texttt{.vbn} à la fin)
du paramètre qui contient la datation de cette
variable; ce champ n'est généralement pas présent
dans les variables qui contiennent elle-même un temps, mais il n'est pas
vraiment interdit qu'il y soit présent, dans les cas où l'on
souhaiterais comparer plusieurs systèmes, et donc savoir à
quelle heure (du premier système) correspond l'heure (du deuxième).
   \item[time\_bnds] dans ref\_time, précise les bornes temps de 2 manières possibles : 
   
time\_bnds = vector : <nom du vecteur de bornes temps> 

time\_bnds = decal\_av, decal\_ap (exemple : time\_bnds = -1800, 0 ) 

les bornes temps sont calculées par time+decal\_av, time+decal\_ap 

le nom du vecteur bornes temps est sous la forme : rft\_bnds\_<chaine libre>\_freq
    \item[datation\_convention] attribut informatif dans le reftime pour indiquer la norme de codage
     
	datation\_convention = N avec N variant entre 0 et 1  (0, : debut, 1 : fin intervalle)
    \item[type\_ech] Indique la façon dont la variable a
   été échantillonnée. Deux valeurs sont 
  possibles : \textsf{regulier} et \textsf{irregulier}.
   Ce champ n'est obligatoire que pour les vecteurs de datation; 
  S'il n'est pas présent dans un vecteur de données, on le déduit du
vecteur de datation (celui désigné par
\textsf{ref\_time}); et si celui-ci est absent, on finit par supposer que
l`échantillonnage est irrégulier.
    \item[freq\_ech] Si la variable a été
régulièrement échantillonnée, la valeur de la
fréquence d'échantillonnage, en Hz\footnote{Attention, cette
fréquence n'est pas forcément un entier $\ge$ 1, c'est un
réel $>$ 0}. Cette valeur peut être écrite avec toutes les décimales
nécessaires (dans une limite raisonnable, disons 30 caractères), mais
les expressions mathématiques (par exemple $ 1 / 3$) ne sont pas admises.
Pour les variables irrégulièrement échantillonnées (\textsf{type\_ech = irregulier}), ce champ est
présent et vaut 0.\\
  (Pour préserver une compatibilité
avec les précédentes versions du format, qui utilisaient le
mot-clef \textsf{sample}, les programmes qui lisent des fichiers vecteurs
doivent prévoir le cas où \textsf{freq\_ech} est absent, et
lire alors le champ \textsf{sample}).
    \item[unit] une unité répertoriée par la librairie \textit{udunits} du
                groupe unidata ou \textsf{unitless} (pour les valeurs issues
                d'un système de numérisation, l'unité la plus appropriée est
                \textsf{bit}). \\
                Pour des raisons historiques, on peut trouver l'unité
                \textsf{without} qu'il convient alors d'interpréter comme
                \textsf{unitless}.
%%%%On trouvera en annexe 1 une liste
%%%%des unités les plus fréquemment employées.
  \item[orientation] lorsque le paramètre contenu dans un vecteur est
une grandeur physique orientée (un angle, une distance sur un axe, une
composante d'un vecteur), ce champ contient une description de cette
orientation. Par exemple, pour un roulis : «~>0 lorsque le coté droit
de l'avion (ou du navire) est enfoncé~». La présence de cet attribut n'empêche
pas d'utiliser un nom explicite, lorsque c'est possible, par exemple
"\texttt{ray\_glo\_d\_...}" pour un rayonnement global descendant, ou ``\texttt{vit\_e\_\ldots}"
pour une composante de vitesse positive vers l'est.

  \item[date\_creation] au format \textsf{jj/mm/aaaa hh\string:mm\string:ss}.
  \item[origine] est un champ qui décrit à partir de quoi la variable a été
    crée.  Il est composé de plusieurs mots séparés par des espaces. Le
    premier décrit la fonction du programme producteur : c'est soit
    \textsf{lecture}, soit \textsf{calibration}, soit \textsf{synchronisation}, 
    soit enfin \textsf{traitement}.\\

    La suite du champs précise la description. Elle varie selon les cas :

    \begin{description}
      \item{Dans le cas \textsf{lecture}}, le mot suivant est le type de
        fichier d'origine [bande/f\_turb/NetCDF,...] puis le nom (sans le
        chemin complet) si il s'agit d'un fichier.

      \item{Dans le cas \textsf{calibration}}, le mot suivant est le nom du
        fichier vecteur d'origine, suivi de l'éventuelle méthode employée pour
        la calibration.

      \item{Dans le cas \textsf{traitement}}, la suite indique le  nom du
        programme de traitement (par compatibilité historique, si la suite est
        vide, on considérera que le programme d'origine est adc).
    \end{description}

\end{description}

%%%\paragraph{attributs optionnels}
  Les champs suivants sont optionnels, mais il est recommandé de les
renseigner chaque fois que c'est possible :
\begin{description}
   \item[name] Nom complet du paramètre, ne dépassant pas une
ligne, en français.
   \item[english\_name] Nom complet du paramètre, ne dépassant 
pas une ligne, en anglais (dans un fichier NetCDF, ce serait l'attribut
\textsf{long\_name}).
   \item[information] Description textuelle élargie, pouvant compléter 
le champ «name». 
  \item[date\_vol] la date de mise en route de l'acquisition, au 
format \textsf{jj/mm/aa hh\string:mm\string:ss}. L'information
sur le jour doit correspondre au champ date de la section [IDENTIFICATION]
du fichier .inf de l'acquisition. Si ces informations sont contradictoires,
le fichier .inf fait foi.
%% Ajouts suite à la réunion CAM/TRAMM du 02/12/1999
   \item[bornes] Dans le cas où chaque échantillon du paramètre est un 
tableau (attribut \textsf{dim\_array} supérieur à 1), ce champ précise 
les bornes de chacun des éléments du tableau.
Les bornes peuvent être être définie littéralement, sous la forme d'une expression permettant de les
calculer, ou être contenues dans un autre vecteur pour le cas où les bornes sont variables dans le temps.
La syntaxe peut prendre 3 formes :
\begin{itemize}
\item une suite de couples sous la forme :\\
\texttt{bornes = ($x_1$,$y_1$); ...; ($x_n$,$y_n$) }.\\
Les $x_i$ sont en général des réels, ou le caractère - si la
limite n'est pas définie (+ ou - l'infini). Le nombre de
couples est la valeur de \textsf{dim\_array}. (\textbf{Exception} : Les
vecteurs "quanta" provenant de la carte-coupleur \textsc{CNRM-1D} contiennent
une structure complexe : tableau de comptage et voies annexes. Dans 
ce cas, le champ \textsf{bornes} ne décrit que le tableau de comptage,
et le nombre de couples correspond à la valeur de l'attribut \textsf{classes}).

\item une formule permettant de calculer les bornes : \\
bornes = <range> : <debut> $d$, <fin>$f$, <nclasses>$N$, <pas>$P$, <largeur>$L$
Les valeurs des champs $d$, $f$, $N$, $P$, $L$ sont des valeurs numériques réelles.\\
L'ordre des champs n'a pas d'importance et 3 d'entre eux sont obligatoires.\\
Le champ <pas> est obligatoire dans tous les cas.\\
Le champ <largeur> est optionnel.  S'il n’est pas présent, alors $L = P$ \\
Parmi les 3 champs : <debut>, <fin> , <nclasses>, 2 d’entre eux sont obligatoires \\
La relation entre les champs est :  $f = d + (N -1)*P + L$ 

\item la réference à un  autre vecteur contenant les bornes variables dans le temps : \\
<vector> : \textsf{vecteur\_bornes}. Ce vecteur a 2 dimensions \textsf{dim\_array} et 2. Il contient à chaque pas de temps les bornes de début et fin de chaque classe.\\ 

\end{itemize}

Pour les vecteurs à plusieurs dimensions, les bornes de chaque dimension sont décrites par 
les mêmes syntaxes séparées par des \textsf{|} si plusieurs dimensions. \\
Par exemple : 

\-\hspace{1cm} {bornes = <range>: .... | <vector>: ...  | ($x_1$,$y_1$); ...; ($x_n$,$y_n$)}

L'ordre suit la même règle que \textsf{dim\_array}.

   \item[unit\_bornes] Indique l'unité de l'attribut \textsf{bornes}
décrit ci-dessus, séparés par des \textsf{|} s'il y a plusieurs dimensions.
   \item[type\_acq]indique la façon dont la variable a
    été acquise. Deux valeurs sont possibles :\\
    \textsf{synchrone} (dans le cas où l'instant de l'acquisition est
     commandé par le système d'acquisition),\\
     ou \textsf{asynchrone}
      (dans le cas où l'appareil lui-même décide du moment où
      il envoie ses données au système d'acquisition).
   \item[position] Un mot qui désigne un emplacement sur la plate-forme
    de mesure.

   \item[misval] la valeur spéciale pour les cas où la valeur
           n'est pas disponible.
   \item[invval] la valeur spéciale pour les cas où la valeur
           était manifestement hors-norme.
   \item[obs\_min] La plus petite valeur contenue dans la série
   \item[obs\_max] La plus grande valeur contenue dans la série
   \item[normal\_min] La plus petite valeur physiquement possible.
   \item[normal\_max] La plus grande valeur physiquement possible.
   \item[famille] Le type de grandeur physique de la mesure. La liste des
familles connues est contenue dans le fichier \path{famille}, décrit
en section \ref{sec:familles}, page \pageref{sec:familles}.
   \item[fct\_origine] Une description de la fonction de calcul qui a
produit ce vecteur.  Dans les cas où l'origine est «~lecture~» ou
«~calibration~», il peut servir à donner des précisions sur le calcul
effectué (formule, valeur des coefficients). Dans le cas où l'origine
est «~traitement~», il doit indiquer
tous les éléments d'entrée du calcul (vecteurs, coefficients), ainsi
que les éléments de sortie (vecteurs ou paramètres produits
simultanément). Dans le cas de traitement par adc, ce champ est une copie
de la ligne de commande du calcul.
   \item[fct\_comment] Commentaires ou compte-rendu produit(s) par la
fonction de traitement.
   \item[origine\-WARNING] messages d'avertissement ou de détection
d'anomalie qui ont été émis lors de la fabrication du fichier-vecteur
(par exemple « unité non reconnue »).

\end{description}

  Pour les vecteurs issus directement d'instruments (les \og quanta\fg{}
et \og cal\fg{} ), il existe deux méthodes (deux séries de champs optionnels)
pour décrire la mesure. Les programmes recherchant cette information doivent
donc savoir traiter les deux cas.

Le premier modèle est simple, et se base sur trois attributs :
\begin{description}
   \item[sensor] : nom de l'appareil de mesure.
%% ajouts 2005
La désignation complète d'un instrument comporte : le nom du constructeur,
le nom du modèle, le numéro de série, et une éventuelle référence interne 
(inventaire, ...).
Pour un meilleur traitement automatique ultérieur, il est recommandé de  
fournir ces quatre informations, séparées par des ":".
%% Fin ajouts 2005
   \item[transducer] : transducteur. Format exact non spécifié.
   \item[calibration] : méthode et date de calibrage ou étalonnage. Format exact
non spécifié.
\end{description}

Le second modèle est plus complexe, mais plus détaillé :
\begin{description}
%% Fin ajouts suite à la réunion CAM/TRAMM du 02/12/1999
   \item[chaine\_acquisition] (anciennement chaine\_mesure obsolète) : pour les paramètres bruts, description de la chaîne 
d'acquisition, selon le format : \\
\texttt{chaine\_acquisition = module\_1 module\_2 module\_3 module\_4 ...}

Les noms des modules sont de la forme module\_N, toujours en minuscules,
et séparés par des espaces.

   \item[module\_N] : description du N\up{ieme} module de la chaîne
     d'acquisition, suivant le format : \\
     \texttt{module\_N = <type\_module>:<fabricant>:<Modèle>:<Numéro de série>:<Numéro de référence interne>:<date et lieu de calibrage ou étalonnage>}

\smallbreak
Il n'est pas interdit qu'il y ait plusieurs modules de même
type dans une chaîne d'acquisition.\\

\smallbreak
La définition précise des types est : \\
  \begin{tabular}{rp{10cm}}
    antenna & Élément mécanique permettant de capter
    soit le flux d'air et de l'amener sur l'élément sensible,
    soit un signal radio-électrique (radars, récepteurs GPS, ...) \\
    
    sensor & Élément sensible. Une de ses propriétés physique varie en fonction
    de la grandeur que l'on veut mesurer. \\
    
    transducer & Transforme une propriété physique en tension, fréquence, 
    courant, ... Il peut être constitué d'un ensemble sensor+transducer 
    lorsqu'ils sont fortement liés (sensor pas facilement interchangeable) \\
    
    conditionner & Circuit électronique de traitement du signal issu 
    d'un capteur ou d'un transducteur. Il est utilisé pour simplifier l'interface
    du capteur avec le système d'acquisition de données. Il intervient entre le 
    capteur ou le transducteur, et l'unité d'acquisition. \\
    
    acquisition\_unit & Interface entre le monde extérieur (tensions, bus 
    numériques, lignes RS232) et le système d'acquisition.\\
    
    instrument & Appareil autonome se suffisant à lui-même, disposant 
    d'un affichage ou d'un système de stockage des données.\\
  \end{tabular}
\end{description}

%%%\paragraph{attributs libres}
  Enfin, tous les noms qui ne sont pas décrits ici sont libres d'utilisation
pour certaines applications spécifiques. En particulier les vecteurs résultant
d'une simple lecture, ou d'une calibration, conservent, pour que
l'auto-documentation soit complète, les attributs qui sont définis dans le
fichier d'origine (en-tête de fichier NetCDF, fiches de configuration du
système d'acquisition du CAM, etc...).

  Le tableau suivant recense les mots déjà définis par le système
d'acquisition du CAM, que d'autres applications doivent donc éviter
d'utiliser : \\

  \begin{tabular}{p{5cm}p{5cm}p{5cm}}
      bloc                 & calibration                 & channel      \\
      classes\footnotemark[1]
                           & cnv\texttt{*} (cnvi, cnvj, cnvrtn, ...) & codage \\
      code                 & codeRT                      & coupleur     \\
      cutoff\_freq         & delay                       & gain         \\
      index                & label                       & law          \\
      laws                 & mode                        & module       \\
      msgindex             & range\footnotemark[2]
                                                         & raz          \\
      recup                & reftime\footnotemark[3]     & resolution   \\
      sample\footnotemark[4]
                           & transducer                  & typeofdata   \\
      version              & voiesup\footnotemark[1]     & voie\_sup\_\textit{x}\\
  \end{tabular}

%%% notes de bas de page
\footnotetext[1]{\textsf{classes} et \textsf{voiesup} servent à préciser
  \textsf{dim\_array} pour les vecteurs «quanta» de sondes 1D. Normalement,
  \textsf{dim\_array}=\textsf{classes} + \textsf{voiesup}.}
\footnotetext[2]{Attention : cet attribut n'a pas le même sens pour les données
  issues de l'INS et des cartes analogiques.}
\footnotetext[3]{Attention : ne pas confondre avec ref\_time !!!}
\footnotetext[4]{Attention : n'est pas obligatoirement exprimée en Hz ! Voir
  aussi le mot-clef «blocage» dans les descripteurs de blocs}  

\subsection{Format de la partie binaire}
Cette partie du fichier contient la suite des
échantillons de la variable, dans le format indiqué dans la
partie descriptive.

    Le cas des variables qui ne sont pas de simples scalaires, mais des
tableaux (\texttt{dim\_array>1}) nécessite quelques précisions.  Le principe
général est que chaque tableau est considéré comme un échantillon, et le fichier
contient donc la suite chronologique des tableaux. Lorsqu'un tel fichier sera
chargé directement en mémoire, la façon de le décrire sera donc différente
selon le langage informatique utilisé.

Par exemple, pour des données se présentant sous la forme d'un tableau 
unidimensionnel (dim\_array=n), la déclaration correspondante en C sera :\\
\texttt{float data [nombre\_de\_secondes] [taille\_du\_tableau]}, ou encore\\
\texttt{float data [] [taille\_du\_tableau],}

alors qu'en FORTRAN, on écrirait plutôt :\\
\texttt{DIMENSION data (taille\_du\_tableau, nombre\_de\_secondes)}, ou encore\\
\texttt{DIMENSION data (taille\_du\_tableau, *)}

Une autre façon de voir les choses est de considérer un programme "universel", qui 
doit être capable de traiter des variables dont le dimensionnement est différent;
il faut alors considérer la zone de données comme un tableau unidimensionnel, 
et faire lui-même les calculs d'indice. La formule qui donne l'indice $k$ de la $i$\up{ème}
valeur de la $s$\up{ème} seconde de donnée est : $k = i + s*\textrm{taille\_du\_tableau}$.

Les cas de nombre de dimensions supérieures (par exemple les images pour lesquelles
\textsf{dim\_array=m|n}) se traitent de la même façon;
la formule générale qui donne l'indice $k$ du pixel $i$,$j$ de
l'image de la $s$\up{ème} seconde est :
$ k = s*m*n + i*m + j $

\subsection{Cas particulier des variables temps}

Elles contiennent toutes un temps depuis 00h00 UTC le jour de la mise en route du
système d'acquisition, exprimé 
en secondes. (Attention, pour les acquisitions qui chevauchent minuit
UTC, le nombre de seconde peut donc dépasser 86400)

Elles sont toutes au format \textsf{Ieeefloat\_64}, avec
\textsf{dimension = 8}, et \textsf{unit = second}
(le système unidata permettrait de rajouter  since 00h00 ).

Pour éviter de lire tout le fichier lorsque l'on recherche juste les heures de
début et de fin, on rajoute des champs \textsf{heure\_deb} et \textsf{heure\_fin},
au format \textsf{hh\string:mm\string:ss.ddd}

\section{Nommage}

  L'extension de nom utilisée pour ces fichiers est : \texttt{.vbn}.

  Le corps du nom de fichier (``\texttt{basename}'') correspond au
nom de la variable, tel qu'il est défini dans l'en-tête.

  Le nom d'un vecteur est à la fois un nom de fichier, et un nom de
variable, qui est utilisé pour faire référence à ces données (sélection
pour un graphique, formule de calcul, ...). La liste des caractères
possibles dans ce nom
cumule donc les contraintes des systèmes de fichiers et des interpréteurs.

\medbreak
  Les caractères suivants sont explicitement autorisés :

\begin{itemize}
	\item les lettres : de A à Z, en majuscule et en minuscul.
	\item les chiffres :  de 0 à 9
    \item le caractère souligné : « \_ »
    \item le caractère pourcent : « \% »
    \item le caractère 'deux points' : « : »
    \item le caractère 'deux points' : « : »
    \item le caractère arobase       : « @ »
    \item le caractère point         : « . »
\end{itemize}

\medbreak
  Il peut être envisagé d'utiliser les caractères suivants, qui ne sont pas
explicitement interdits  :
%– Enfin, il est recommandé que les programmes sachent manipuler des noms
%contenant les caractères qui n'ont pas encore été mentionnés («!#,:@[]
%{} ») mais le support de ces caractères n'est pas une obligation.

\begin{itemize}
  \item le caractère 'point d’exclamation' : « ! »
  \item le caractère 'dièze'               : « \# »
  \item le caractère 'virgule'             : « , »
  \item les caractères accolades ouvrante et fermante : « \{ » et « \} »
\end{itemize}

\medbreak
  A titre documentaire, les raisons d'exclusion des autres caractères sont
données en annexe, page \pageref{sec:exclu}

  D'un point de vue sécurité, l'interdiction de certains caractères ne veut
pas dire que les scripts ne doivent pas se protéger contre les noms de fichiers
qui contiendraient ces caractères. Cela signifie seulement qu'il est légitime,
dans le cadre AIDA, d'ignorer ou de signaler ces fichiers.

\textbf{Remarque} : les règles de nommage des paramètres (règles de
construction des noms en fonction de taxonomies et référentiels) ne sont pas
purement techniques, et sont donc décrites dans le volume 2 : manuel
d'utilisation (\url{user.pdf}).

